\section{Введение и постановка}

\paragraph{Задача.} Нужно построить последовательную рекомендательную модель
на реальных данных прослушиваний, снять метрики ранжирования на официальном
протоколе, обучить reward model на сигналах поведения и дообучить политику
методами обучения с подкреплением, исследовав, как меняется поведение выдачи.

\paragraph{Формулировка.} Задачу мы решаем как офлайн-оптимизацию политики.
Состояние --- история прослушиваний пользователя, действие --- айтем,
награда --- оценка reward model для пары (история, айтем), переход ---
детерминированное добавление айтема в историю. Обучение с подкреплением нигде
не взаимодействует с реальными пользователями, поэтому основным предметом
исследования становится поведение награды и политики под оптимизацией, а не
эффект в проде.

\paragraph{Гипотезы.} Гипотезы сформулированы так, чтобы каждую можно было
опровергнуть одним числом (таблица~\ref{tab:hyp}).

\begin{table}[t]
\centering
\caption{Гипотезы, критерии и итог. Столбец «итог» заполнен по результатам
полного прогона (раздел~\ref{sec:results}).}
\label{tab:hyp}
\small
\begin{tabularx}{\textwidth}{@{}clLc@{}}
\toprule
\# & Гипотеза & Критерий & Итог \\
\midrule
H1 & RL повышает награду RM & рост награды при сдвиге KL & да (+4.4\%) \\
H2 & За награду платят релевантностью & $\Delta$NDCG@10 $< 0$ вне 95\% CI & опровергнута \\
H3 & Сдвиг к «дешёвым» айтемам & длина $\downarrow$, Jaccard $\uparrow$ & да \\
H4 & Многошаговая генерация деградирует & падение уникальности, рост повторов & да \\
H5 & Долгосрочная награда ведёт себя иначе & сравнение с R1 & за рамками работы \\
H6 & Коллапс ограничен якорем KL & связь $|\Delta\text{NDCG}|$ и $\beta$ & $\beta$ фиксирован \\
H7 & Награда и релевантность расходятся & знак изменения по головам & да \\
\bottomrule
\end{tabularx}
\end{table}

H2 проверена: улучшения нет, падение выходит далеко за пределы шума и идёт
в одну сторону у большинства затронутых пользователей.
H5 требует симулятора с обратной связью и остаётся за рамками работы;
H6 в этой постановке сводится к константе: $\beta_{\text{KL}}$ фиксирован,
зависимость эффекта от якоря измеряется программой абляций
(раздел~\ref{sec:limits}). Критерий H7 выполняется: голова «досмотр»,
которую оптимизирует награда, растёт, а голова «лайк» падает.

\paragraph{Вклад.}
\begin{enumerate}[leftmargin=1.2em,itemsep=1pt,topsep=2pt]
  \item воспроизводимый ETL на официальном протоколе Yambda
        (300 дней обучения, 30 минут разрыва, день валидации, день теста);
  \item измеренная цена обучения с подкреплением: награда $+4.4\%$,
        NDCG@10 $-84\%$ при парном бутстрэпе по 4\,600 пользователям;
  \item три независимых признака коллапса поведения: падение длины трека,
        рост межпользовательского Jaccard и расхождение знаков у голов
        reward model;
  \item публичный код, в котором каждая цифра получается одной командой
        через стадии~\code{run.py}.
\end{enumerate}
